% Generated draft: requires mathematical and source review.
\documentclass[11pt]{article}
\usepackage[margin=25mm]{geometry}
\usepackage{fontspec}
\setmainfont{Latin Modern Roman}
\usepackage{amsmath,amssymb,mathtools,unicode-math}
\setmathfont{Latin Modern Math}
\usepackage{enumitem,needspace,xurl,hyperref}
\setlength{\parindent}{0pt}
\setlength{\parskip}{5pt}
\setlength{\emergencystretch}{3em}
\newcommand{\N}{\mathbb{N}}
\newcommand{\Z}{\mathbb{Z}}
\newcommand{\Q}{\mathbb{Q}}
\newcommand{\R}{\mathbb{R}}
\newcommand{\C}{\mathbb{C}}
\newcommand{\F}{\mathbb{F}}
\newcommand{\A}{\mathbb{A}}
\newcommand{\PP}{\mathbb{P}}
\newcommand{\E}{\mathbb{E}}
\newcommand{\eps}{\varepsilon}
\newcommand{\dd}{\,\mathrm{d}}
\begin{document}
\section*{UnsolvedMath ID: 30001623}
\textbf{Source problem code:} catalogue entry
\subsection{Definitions and mathematical statement}
The supplied record concerns one-dimensional hierarchical coalescence models, but does not define their state spaces, initial laws, coalescence rules, or hierarchy parameters.

A scaling limit is a limiting law of suitably rescaled random quantities associated with a model. To describe one possible formulation precisely, suppose that rescaled observables have probability laws \(\mu_{M,n}\) on a common metric space \(S\), where \(M\) denotes a model and \(n\) indexes its hierarchy. Weak convergence to a probability law \(\mu_M\) means
\[
\mu_{M,n}\Rightarrow\mu_M
\quad\Longleftrightarrow\quad
\int_S f\,\dd\mu_{M,n}\longrightarrow\int_S f\,\dd\mu_M
\quad\text{for every bounded continuous }f:S\to\R.
\]
In such a formulation, universality would mean that suitably normalized observables from different models in an identified class have the same limiting law. This is only an explanatory schema: the source does not specify these observables, normalizations, model class, hierarchy parameter, state space, or mode of convergence, and does not establish that this particular schema is the intended formulation.

\textbf{Source question.} How can one explain the existence of scaling limits, and why do very different one-dimensional hierarchical coalescence models exhibit the same limiting distributions?

The requested target is an explanation of scaling-limit existence and shared limiting distributions across such models. The supplied record does not give hypotheses or quantifiers sufficient to turn this question into a uniquely specified theorem or conjecture. In particular, it does not assert that every one-dimensional hierarchical coalescence model has a scaling limit or shares one universal law.
\subsection{Short English statement}
Explain why appropriate rescalings of certain one-dimensional hierarchical coalescence models converge in distribution, and why different models can share limiting distributions. The source leaves the models, initial conditions, rescalings, observables, and convergence notion unspecified, so this remains a research question requiring mathematical clarification.
\subsection{Sources}
\begin{itemize}
\item Supplied local UnsolvedMath registry record, permanent ID \(30001623\), title \emph{Universal Scaling Limits in Hierarchical Coalescence}, source problem code \texttt{catalogue entry}; collection: Oberwolfach Reports Open Problems. The project registry is \texttt{data/unsolvedmath/problems.json}, with additional existing records in \texttt{catalog/problems.json}; the supplied record's location within these files has not been independently checked.
\item Supplied report citation: \emph{Large Scale Stochastic Dynamics (2011)}, Oberwolfach Reports, DOI \texttt{10.4171/owr/2010/50}, pp.\ 2877--2975. Supplied source URL and literature source labeled ``original OWR report'': \url{https://doi.org/10.4171/owr/2010/50}. The record attributes the question to workshop open-problem discussions; no individual proposer is supplied. The proposed year is recorded as 2010.
\item Literature assessment dated \texttt{2026-08-21}: the supplied curation labels the problem open and states that selected model-specific limits have been proved, while the general explanation of scaling limits and universality remains a broad research question. This is supplied dated curation, not an independently verified assessment. The supplied background expressly requests independent verification of time-sensitive status and cited claims.
\item The registry labels the statement ``exact'' and says it was retained byte-for-byte because it is internally readable. These are supplied record qualifications; they do not constitute independent verification against the report.
\end{itemize}
\end{document}
